\documentclass[../../main.tex]{subfiles}
\graphicspath{{\subfix{../../image/}}}

\begin{document}

\section{参考文献}\label{section:ch02-references}

\textbf{\S2.2.} 关于Laplace方程与 Poisson 方程的进一步内容，Gilbarg--Trudinger \cite{G-T} 第 2--4 章是很好的资料。解析性的证明取自 Mikhailov \cite{M}。J. Cooper 在 Green 函数方面给予了我帮助。

\textbf{\S2.3.} 关于热方程的进一步信息，见 John \cite{J2} 第 7 章或 Friedman \cite{Fr1}。\cref{theorem:ch02-heat-mean} 归功于 N. Watson（Proc. London Math. Society 26 (1973), 385--417），\cref{theorem:ch02-heat-strong-max} 的证明亦然。\cref{theorem:ch02-heat-cauchy-max} 取自 John \cite{J2}，\cref{theorem:ch02-heat-smoothness} 沿 Mikhailov \cite{M}。\cref{theorem:ch02-heat-backs-unique} 取自 Payne \cite{Pa} \S2.3。

\textbf{\S2.4.} 见 Antman（Amer. Math. Monthly 87 (1980), 359--370）对一维波动方程作为振动弦模型的细致推导。此处给出的波动方程解法沿 Folland \cite{F1}、Strauss \cite{St2}。

\textbf{\S2.5.} J. Goldstein 贡献了习题 24。

\end{document}